% Overrides Quarto's own before-body.tex (just a conditional \maketitle --
% unused, since documents here build their own title page by hand).

% A document that fits on one page (cv-blocks.compact: true) needs no
% running header/footer -- switch off fancyhdr's pagestyle for it.
% Independent of cv-blocks.long: that one selects long_only/short_only
% *content*, not page count -- a document can show the short content set
% and still span several pages (e.g. a dossier scientifique with
% long: false), so it must keep the fancy pagestyle from
% include-in-header.tex by default.
$if(cv-blocks.compact)$
\pagestyle{empty}
$endif$

% Quarto's own template unconditionally loads `parskip` (or falls back to
% \setlength{\parskip}{6pt plus 2pt minus 1pt}), adding visible vertical
% space between consecutive \multblock/\block entries -- each one is
% emitted as its own paragraph (blank-line-separated) by cv-patterns.lua,
% so every entry boundary is a new paragraph break under parskip's rules.
% Reset to LaTeX's classic zero-\parskip so entries stay as tightly
% spaced as the original cv.cls (which never loaded parskip).
\setlength{\parskip}{0pt}

% \nom/\title (used in the fancyhdr footer below, and freely in a
% document's own raw LaTeX, e.g. a hand-built title page) come from
% Quarto's native author/title metadata -- see the extension README,
% "Usage". Every consuming document is expected to set both. `\title` is
% already a real article.cls command (its argument sets \@title, used by
% \maketitle, which this extension never calls) -- \renewcommand
% repurposes it as a plain text-holding macro instead.
\renewcommand{\title}{$title$}
\def\nom{$for(authors)$$it.name.literal$$sep$ \and $endfor$}

% Used by the fancyhdr footer in include-in-header.tex.
\newcommand{\footer}{\nom\ -- \title}

% Colors. Referenced by `linkcolor`/`urlcolor`/`citecolor` (native
% metadata, resolved into \hypersetup{...} in the preamble) and by
% \titleformat below -- works even though this runs after \begin{document}
% since hyperref/xcolor resolve color names lazily, only when typeset.
\definecolor{mblue}{rgb}{.204,.353,.541}
\colorlet{mred}{red!80!black}

% Small caps via URW Avant Garde instead of LaTeX's default (faked, from
% the current font) small caps. \fontencoding{T1} is needed under
% LuaLaTeX, which defaults to TU encoding and only has T1 shapes for
% "pag" -- harmless under pdflatex too, since T1 is already the active
% encoding there (include-in-header.tex loads fontenc[T1] unconditionally).
\renewcommand{\sc}{\fontencoding{T1}\fontfamily{pag}\fontseries{m}\fontshape{sc}\selectfont}
\renewcommand{\textsc}[1]{{\sc#1}}

\titleformat{\part}{\sc\huge\color{mred}}{}{0pt}{\noindent\hrulefill~#1~\hrulefill}
\titleformat{\section}{\sc\huge\color{mred}}{}{0pt}{\noindent\hrulefill~#1~\hrulefill}
\titleformat{\subsection}{\sc\Large\color{mred}}{}{0pt}{\noindent#1~\hrulefill}
\titleformat{\subsubsection}{\sc\large}{}{0pt}{\noindent#1}
\titleformat{\paragraph}[runin]{\sc}{}{}{#1}[]
\titleformat{\subparagraph}[block]{\sc}{}{}{\noindent#1}[]
\titlespacing{\subparagraph}{0pt}{0pt}{0pt}

% \multblock/\block: the two primitives cv-patterns.lua's raw LaTeX
% output is built from -- see the extension README ("How it works").
\newcolumntype{R}[1]{>{\raggedleft\let\newline\\\arraybackslash\hspace{0pt}}m{#1}}

\newcommand*{\multblock}[3]{%
  \par\noindent\begin{tabular}[t]{@{}R{.15\linewidth}|p{.8\linewidth}}
    \hfill{\sf#1} & {\sc#2} \\\relax
    \readBlock#3\relax\relax
  \end{tabular}\medskip
}

\newcommand*{\block}[2]{%
  \par\noindent\begin{tabular}{@{}p{.15\linewidth}|p{.8\linewidth}}
    {\hfill\sf#1} & {\sc#2} \\
  \end{tabular}\medskip
}

\newcommand{\readBlock}[2]{%
  \ifx\relax#1\empty
  \else
  \hfill{\em#1} & #2\\
  \expandafter\readBlock
  \fi
}

% Profondeur 1 de la table des matières
\setcounter{tocdepth}{3}
